\section{Introduction}

Ability-based epistemic logics arise in many areas of computer science, where they provide a formal framework for strategic reasoning and AI planning. Several logics in this family have been studied in recent years. In particular, epistemic logics formalising the notion of \textit{knowing how} have received considerable attention because of their conceptual simplicity and strong meta-logical properties.\par
The simplest version of \textit{knowing how}, $\mathcal{L}_{Kh}$, based on the paper by (Wang 2015\cite{Wang2015}), with the semantics reformulated in (Areces et al. 2021\cite{Areces2021}). 
This logic captures the agent's ability to guarantee the completion of a task under a given condition. 
This idea is ideal in a sense, since even though the agent manages to do something under a condition, it does not imply the agent really know how. 
For example, one can win a lottery does not imply one can guarantee how to win. The uncertainty is involved in this case. 
Therefore, another logic may be preferred, $\mathcal{L}^U_{Kh}$. In this logic, an agent may have \textit{indistinguishability} among several plans they might take, a critical concept in epistemic logic. 
Indistinguishable plans form equivalence classes, partitioning the set of available plans. Intuitively, plans considered identical by the agent are placed in the same class. 
Furthermore, this logic involves multi-agents, as opposed to a single agent. Finally, a new logic, $reg\text{-}\mathcal{L}^U_{Kh}$, a variation of the second version, 
introduces equivalence classes that are defined in terms of finite automatons. \par
Additionally, for demo and testing purpose, we've implemented a simple web-application to support model-checking by offering a visual experience and keeping a history of formulas and models (within a session). 